% Hardware Specifications for Thesis
% Generated automatically - use in Chapter 4 Methodology

\subsubsection{Specyfikacja sprzętowa}
\label{subsubsec:specyfikacja-sprzetowa}

Wszystkie testy wydajnościowe przeprowadzono na komputerze o następującej specyfikacji:

\begin{itemize}
    \item \textbf{Procesor}: AMD Ryzen 9 7900X3D 12-Core Processor (24 rdzenie, 48 wątków)
    \item \textbf{Karta graficzna}: NVIDIA GeForce RTX 3090
    \item \textbf{Pamięć GPU}: 24 GB GDDR6X
    \item \textbf{Sterowniki NVIDIA}: wersja 590.48.01
    \item \textbf{Pamięć RAM}: 32 GB
    \item \textbf{System operacyjny}: Arch Linux (jądro Linux 6.18.5-arch1-1)
    \item \textbf{Dysk}: SSD o pojemności 3,6 TB
\end{itemize}

\subsubsection{Specyfikacja oprogramowania}
\label{subsubsec:specyfikacja-oprogramowania}

W badaniach wykorzystano następujące wersje oprogramowania:

\begin{itemize}
    \item \textbf{Unity}: 6.0 (6000.0.58f2) LTS
    \item \textbf{Unreal Engine}: 5.5.3
    \item \textbf{NVIDIA Nsight Systems}: 2025.5.2
\end{itemize}

Wybór wersji LTS silnika Unity podyktowany był stabilnością oraz długoterminowym wsparciem, co jest istotne z punktu widzenia powtarzalności badań. W przypadku Unreal Engine wybrano najnowszą dostępną wersję stabilną w momencie rozpoczęcia badań.

\subsubsection{Konfiguracja testowa}
\label{subsubsec:konfiguracja-testowa}

Przed wykonaniem każdego testu zastosowano następującą procedurę przygotowawczą:
\begin{enumerate}
    \item Zamknięto wszystkie aplikacje działające w tle
    \item Wyłączono automatyczne aktualizacje systemu
    \item Ustawiono tryb wydajności zasilania (performance mode)
    \item Odczekano 5 minut na stabilizację termiczną systemu
    \item Wykonano 3 pomiary dla każdego scenariusza testowego, z których obliczono wartość średnią
\end{enumerate}
